\section{The mixed marked compiler on the restarted history}
\label{sec:actual-mixed-marked-compiler}

We give the price calculation which turns the actual three-alphabet
expansion of Section~\ref{sec:allarity-planted-compiler} into the upper
and lower marked transitions. All marked integrals use the genuine fine
representation $P$. Its auxiliary representation $P^0$ has only its
additive remainder removed. The body packets in the two calculations
are nevertheless the same corrected packets. This distinction is
necessary for the normalization argument below.

\subsection{Support currency and the three letter families}

Fix a fine-to-coarse block map of positive integer side $L$, fine and
coarse budget radii $B_f,B_c$ that are positive integers, and a finite
marker set $W$. Let
$\theta:W\to\Lambda_f$ and $u:W\to\Lambda_c$ satisfy
$u=\operatorname{parent}\circ\theta$. Use the cover charge
\[
 \mathcal D_B(t,X)=\exp(t\,\operatorname{cov}_B X),\qquad
 v=e+5,\quad e\ge0,
\]
where $\operatorname{cov}_B X$ is the minimum number of radius-$B$
cubes covering $X$ whose union is connected. Write $X^\uparrow$ for the prescribed coarse
neighborhood of the parent image of $X$. The radius comparison and
cover estimate proved in Section~\ref{sec:inherited-activity-prices}
give, when $B_f\le2B_c$,
\[
 L\operatorname{cov}_{B_c}(X^\uparrow)
       \le15^4\bigl(L+10\operatorname{cov}_{B_f}X\bigr).
\]
Consequently, if $10\cdot15^4(v+2)\le L\sigma_f$, then
\begin{equation}\label{eq:mixed-cover-charge}
 \mathcal D_{B_c}(v+2,X^\uparrow)
 \le T_e\mathcal D_{B_f}(\sigma_f,X),
 \qquad T_e=\exp(15^4(v+2)).
\end{equation}
Indeed multiply the cover inequality by $(v+2)/L$ and exponentiate.
This pays the additive covering cost before taking any marker sum.

All nonnegative sum estimates below are first proved on finite
subfamilies and then extended by monotone convergence.
There are three letter families in one fixed sector and tag:
body letters with empty marker support, lifted background packets,
and nonempty-marker insertion letters. Coalesce background packets
only when their complete activity address and marker support agree.
Their zero-marker coefficients must reproduce the fine archived
activity addresswise. Let $z_a\ge0$ be the body prices. For a
background packet with fine support $X_b$ and paid price $A_b$, use
\[
       f_b=A_b/\mathcal D_{B_f}(\sigma_f,X_b).
\]
The input global packet bound is valid on every fine field, so evaluating
it on the actual chart gives the pointwise bound $f_b$, including on
unhealthy fields. For a regular insertion the analogous estimate holds
on its healthy support; its original mask makes its value zero elsewhere.
Division by the positive cover charge is therefore valid in both cases.
Positive marker weights are included in these prices throughout.

Assume the background site sums obey
$\sum_{b:x\in X_b}A_b\le\eta$. Equation~\ref{eq:mixed-cover-charge}
implies, for each marker $w$,
\begin{equation}\label{eq:mixed-background-marker}
 \sum_{b:w\in F_b}f_b\mathcal D_{B_c}(v+2,X_b^\uparrow)
                         \le T_e\eta.
\end{equation}
Here $\theta(w)\in X_b$ whenever $w\in F_b$. There is no factor
$L^4$ in this marker estimate.

The corresponding coarse site estimate does have that factor. Choose
one genuine fine site in the nonempty supported inventory of every
background packet. At a fixed coarse site at most $L^4$ fine sites
have that parent. The anchor-to-site estimate loses one cover exponent
and has cost
\[
 \mathcal W(B)=(33B)^4\sum_{j\ge0}(j+1)^4e^{-j}.
\]
For completeness, for an anchor $a\in X$ and $x\in X$ the cover
number is at least $\lfloor\operatorname{dist}(a,x)/(8B)\rfloor$.
The number of anchors in its $j$th distance shell is at most
$(33B)^4(j+1)^4$. Summing the exponential loss proves this
period-independent constant. Thus the background site price at
exponent $v+1$ is at most $\mathcal W(B_c)L^4T_e\eta$.

Let $w_h\ge0$ be the insertion-letter prices, and suppose their
marker-root sums at exponent $v+2$ are at most $a$. If every coarse
site receives at most $g$ markers, choose one marker in each letter.
The same argument gives site price $\mathcal W(B_c)ga$ at exponent
$v+1$. On the actual replicated reader geometry $g=n$, the number of
reader labels: the parent cell is fixed by a coarse site, and only its
$n$ labels remain. Combining these bounds with a body site price
$\delta$ gives
\begin{equation}\label{eq:mixed-letter-price}
 \Sigma=\delta+\mathcal W(B_c)L^4T_e\eta
                         +\mathcal W(B_c)ga.
\end{equation}
The total marker-root price at exponent $v+1$ is at most
$T_e\eta+a$. Empty-marker body letters do not contribute to it.

\subsection{Payment of each connected integral}

Fix multipliers $m,l\ge1$. The body-letter product estimate on the
actual sector measure supplies the following stronger integral bound:
if $I$ is a collection of body letters and $X_I$ is the product of
their nonnegative absolute weights, then
\[
 M_I\int X_I\,d\nu\le\prod_{a\in I}z_a.
\]
Here $M_I$ is the product of their individual inflations: a letter
with nonempty forced inventory contributes $m$, and any remaining
body letter contributes $l$. These are the paid unmarked product
estimates, with the same chart and corrected body alphabet as above.
The measure $\nu$ is the Gaussian/Haar probability measure of that
sector and tag. Each other letter has its pointwise price. In addition,
every insertion letter designated as a remainder satisfies one of
\begin{equation}\label{eq:mixed-shadow}
 l h(\xi)\le w_h\quad\hbox{for all }\xi,
 \qquad\hbox{or}\qquad
 h=0\text{ off }E_h,\quad l\nu(E_h)\le1,
\end{equation}
where $0\le h\le w_h$ and $E_h$ is measurable.

Let $C$ be a compatible connected collection of letters. Its full
integrated coefficient, including its marker monomial, is denoted $b_C$.
Then
\begin{equation}\label{eq:mixed-club-integral}
 t_C\|b_C\|_r\le\prod_{j\in C}b_j,
 \quad
 t_C=\begin{cases}
 m,&C\text{ has nonempty forced inventory},\\
 l,&C\text{ is unforced but is not entirely leading},\\
 1,&C\text{ is entirely leading}.
 \end{cases}
\end{equation}
To prove this, split its product into body, background and insertion
factors. The last two products are bounded pointwise by their prices.
If the body set is nonempty, its inflation $M_I$ is at least $t_C$.
In the forced case this uses the exact inventory identity: a forced
site must belong to one of the body inventories. In the unforced
nonleading case any body letter supplies $l$. The paid body integral
therefore proves the estimate.

If the body set is empty, compatibility forces the background set to
be empty as well. A background packet has a nonempty fine inventory,
whereas compatibility identifies the union of those inventories with
the union supplied by the body letters. The component is now a product
of insertion letters. If all are leading, its pointwise bound integrates
against a probability measure. Otherwise choose one remainder letter.
The first alternative in \eqref{eq:mixed-shadow} pays $l$ pointwise;
the second restricts the whole product to $E_h$ and pays $l$ after
integration. This proves \eqref{eq:mixed-club-integral} in every case.
All integrals exist: the body-product majorant is integrable, the other
factors are bounded, and the positive monomial weight can be divided
out. Finite tag averaging preserves the inequality because its weights
are nonnegative and sum to one.

Here is the finite connected-support estimate with the precise smallness
condition used here. Put $c_j=\operatorname{cov}_{B_c}X_j$,
$a_j=b_j e^{vc_j}$ and $\widehat b_j=b_j e^{(v+1)c_j}$, and assume
\[
                (2B_c+1)^4\Sigma\le1.
\]
The sum of $\widehat b_j$ over supports meeting $X_i$ is at most
$|X_i|\Sigma\le c_i$. To count rooted trees, start with
$R_i^{(0)}=a_i$ and iterate
\[
 R_i^{(q+1)}=a_i\exp\left(\sum_{j:X_j\cap X_i\ne\varnothing}
                                       R_j^{(q)}\right).
\]
These are the nonnegative sums over rooted trees of bounded height;
the exponential's factorial divides out permutations of the children.
Allowing repeated labels only enlarges the sum. Induction gives
$R_i^{(q)}\le a_i e^{c_i}=\widehat b_i$ for every $q$.
Each finite connected set containing $i$ has a spanning tree rooted at
$i$, with weight the product of its $a_j$; its distinct labels ensure
that the child permutations cancel exactly the exponential factorials.
Its union charge is no larger than this product charge, by cover
subadditivity. Passing monotonically to the tree limit therefore bounds
the connected-set sum rooted at $i$ by $\widehat b_i$.
A component containing a specified site contains at least one letter
whose support contains that site. Summing the rooted bound gives site
price at most $\Sigma\le e^\Sigma-1$. Likewise a component carrying
a specified marker contains some letter carrying it. Summing those
rooted bounds gives marker-root price at most $T_e\eta+a$. Thus the
same estimates hold simultaneously for every marker, without assigning
independent prices to overlapping subsets. This is a bound on sets of
distinct elementary letters, before the later logarithmic expansion
with repeated labels.
Let the all-hole normalization have raised hole-correction site price
$k$ and boundary-correction site price $\psi$, at exponent $v$; the
raised hole corrections have already been multiplied by $l$. If
\begin{equation}\label{eq:mixed-finish-price}
       e^\Sigma-1+k\le\psi,\qquad T_e\eta+a\le y,
\end{equation}
then the resulting representation has forced price paid by $m$, regular
nonleading price paid by $l$, total regular site price at most $\psi$,
and regular marker-root price at most $y$. Boundary corrections retain
their separately purchased price $\psi$. The corrections carrying no
marker make no contribution to the last bound. This is the complete
price output needed in the marked propagation; the subsequent finite
inventory compression is that of Section~\ref{sec:allarity-marked-propagation}.

\subsection{Two complementary pieces of every leading insertion}

The actual incoming leading letter $W_h$ is split using its one local
cutoff $0\le\Gamma_h\le1$ into $\Gamma_hW_h$ and
$(1-\Gamma_h)W_h$, with the same support and nonempty marker set.
In the squarefree marker algebra their product is zero. Consequently
\[
 (1+\Gamma_hW_hz_{F_h})(1+(1-\Gamma_h)W_hz_{F_h})
                           =1+W_hz_{F_h}.
\]
This is exact also off the healthy set, since $W_h$ retains its original
mask. The first piece is leading; the second is a remainder. The cutoff
plateau implies that the latter vanishes outside the actual warning
event. Its purchased probability is at most $x^{-1}$, where $x$ is
the coarse rung coupling. Thus $l x^{-1}\le1$ supplies the second
alternative in \eqref{eq:mixed-shadow}. Both pieces cost the same
pointwise price, so a leading-family marker budget $h$ costs exactly
$2T_eh$ after coarse support conversion. An arbitrary regular family
of budget $z$ is designated entirely as remainder and costs $lT_ez$;
it uses the first alternative in \eqref{eq:mixed-shadow}.
Hence their union has marker price
\begin{equation}\label{eq:mixed-day-marker}
                    a=lT_ez+2T_eh.
\end{equation}
The event estimate here is the actual joint sector warning estimate
from the local Gaussian construction. A separate Gaussian sample for
each factor would not prove \eqref{eq:mixed-club-integral}.

All leading connected pieces are contained in one cell. Their connected
reader graph has at most $n$ vertices, and the support argument in
Section~\ref{sec:allarity-marked-propagation} bounds its radius by
$(2n+1)\rho$. On the actual elementary support radius $3B_c$, this
is $(6n+3)B_c$. The purchased condition
\[
              6n+3\le A/(3+2n)
\]
therefore puts it inside the required slice radius $AB_c/(3+2n)$.
The whole-mode selection and coefficientwise integration in
\eqref{eq:planted-selection}--\eqref{eq:planted-integral} identify its
projected polynomial with the planted integral on the same cell.
Its integrability also follows directly from the product of its
pointwise letter prices divided by its positive marker weight. Thus
this is the actual leading polynomial, not just a majorant for it.

\subsection{Upper and lower numerical purchases}

Put $H=\exp((\log\log s)^2)$, $K=H^{-1/8}$, and $h=H^{-4}$.
The rung couplings $x,y$ both lie in $[s/2,s]$. Fix $n,e,L$ first.
The large-field allowance $G(y)$ is positive, at most one, and decays
faster than every negative power of $H$ uniformly on this interval.
The radii $B_c=\lceil(\log x)^2\rceil^2$ give
$\mathcal W(B_c)=H^{o(1)}$ uniformly there.

For the upper part use the exponent $e=e_0+2$ and set
\[
 l=\max\left(1,\frac{H^{-3}}{\alpha z+H^{-q}}\right),
 \quad 0\le z\le H^{-4},\qquad m=G(x)^{-2},
 \qquad a=lT_ez+2T_eH^{-4}.
\]
Here $\alpha\ge8T_e$ and $q$ is fixed. If the displayed maximum is
its second entry, $lT_ez\le H^{-3}/8$; if it is one, the same bound
holds eventually since $z\le H^{-4}$. Also $2T_eH^{-4}\le H^{-3}/4$
eventually. Thus $a\le H^{-3}/2$ eventually, after harmlessly
enlarging the threshold. The incoming forced budget is $\eta=G(y)$.
Taking $\delta=K/16$, both nonbody terms in
\eqref{eq:mixed-letter-price} are at most $K/16$ eventually.
Consequently $\Sigma\le3K/16\le K/4$ and
$T_e\eta+a\le H^{-3}$. Also $l\le\max(1,H^{q-3})$, so
$l x^{-1}\le1$ uniformly for $x\ge s/2$ eventually.

For the lower part use $e=e_0$, $l=1$, $m=G(y)^{-1}$,
$0\le z\le H^{-1/2}$ and incoming forced budget
$\eta=z\omega(y)$, where $0\le\omega(y)\le1$ is the actual
lower-interval attenuation. There are no designated leading letters.
The regular marker budget is $a=T_ez$. Again the two nonbody site
terms are each at most $K/16$ eventually. The output marker budget is
\[
 T_ez\omega(y)+T_ez\le2T_ez
       \le\tfrac12\bigl(\alpha z+H^{1/2}z^2\bigr)
                  \quad(\alpha\ge8T_e).
\]
This includes $z=0$ without division. The leading output is empty:
with no body letters, a leading component would consist solely of
insertion letters designated leading, of which there are none; with a
body letter it is already nonleading.

In both cases the hole price is $k=K/4$, the boundary and output
site allowance is $\psi=K$, and the tree smallness condition holds
for sufficiently large $s$. Finally
$e^{K/4}-1+K/4\le K$ for all sufficiently small $K>0$, for example
by $e^t-1\le2t$ on $0\le t\le1$. These calculations prove
\eqref{eq:mixed-finish-price}. The support purchases
\eqref{eq:mixed-cover-charge} for both $e_0,e_0+2$ are made by
choosing $L$ after these fixed exponents and the two positive incoming
support exponents $\sigma_{\rm lower},\sigma_{\rm upper}$. The
separate purchases are
\[
 10\cdot15^4(e_0+7)\le L\sigma_{\rm lower},\qquad
 10\cdot15^4(e_0+9)\le L\sigma_{\rm upper}.
\]
In the actual lower and upper inputs these exponents are respectively
$e_0$ and $e_0+2$. No existence claim is made at a zero incoming
exponent. The coupling threshold is chosen after $L$. For finitely many arities choose the maximum of
their constants before that threshold.

\subsection{The exact representation used by both parts}

The body product and all-hole inputs used above belong to one compiler,
chosen before the marked family. Specifically its body expansion is
\eqref{eq:planted-body} for the genuine $P$, whereas its normalization
is \eqref{eq:planted-holes} for the corrected body alphabet together
with the archived activities of $P^0$. The equality of archived
activities of $P$ and $P^0$ allows that alphabet to embed into the
full marked alphabet, with the same full supports and inventories.
It does not allow the genuine fine exponential to be replaced by that
of $P^0$. The empty-marker component bijection proved in
Section~\ref{sec:allarity-planted-compiler} transports the all-hole
identity to every deleted support of the full marked alphabet.
Applying it at the union of the marked forced supports and using the
exclusion bijection gives precisely the push of the genuine marked
density, normalized by the selected coarse representation.

The unmarked compiler's purchase is made once for multipliers at least
as large as both uses above, and for the larger support exponent.
Decreasing a multiplier or exponent only decreases every nonnegative
price. Its body estimate, hole correction and forced embedding therefore
remain valid for each call. In particular its coarse normalization is
selected before the arbitrary marked background and regular family.
The upper calculation adds the leading family and its complementary
cutoff pieces; the lower calculation uses only the ordinary family.
The same construction supplies both transitions on each rung of the
actual restarted history. The coupling of the push is the genuine
incoming $P.\beta$ in both cases, by the retained course identity.

This proves the mixed integral and price join from the actual local
body and all-hole compiler purchases, the joint warning estimate and
the selected geometry. It does not construct a different restarted
history or deduce those purchases from an abstract Gaussian model.
Their primitive analytic suppliers are the common corrected sector
expansion, its unmarked all-hole normalization, and its paid product
bounds; they must be supplied for the same body alphabet and coarse
representation before this argument is applied.
